\section*{Capítulo \ref{ch:b2:reaction-free-energy} --- Entropía y energía de Gibbs de reacción}

\begin{solution}{exo:b2:reaction-free-energy:1}
(a) Tres moles de gas se convierten en líquido: $\Delta_r S^\circ < 0$. (b) Se forma
un gas: $> 0$. (c) Una molécula de gas da dos: $> 0$. (d) Iones en disolución
forman un cristal: $< 0$. (e) Fusión: $> 0$. (f) $\Delta\nu_{\text{gas}} = 0$:
pequeña, y aquí positiva (\qty{+20}{J/(K.mol)} según las tablas, por ser \ce{HCl}
menos simétrico que \ce{H2} y \ce{Cl2}).
% ledger: s0:HCl_g, s0:H2_g, s0:Cl2_g
\end{solution}

\begin{solution}{exo:b2:reaction-free-energy:2}
$\Delta_r S^\circ = 2(192.77) - 191.61 - 3(130.68) = \qty{-198.1}{J/(K.mol)}$;
$\Delta_r G^\circ = -91.9 - 298.15 \times (-0.1981) = \qty{-32.8}{kJ/mol}$:
\omterm{def:b2:reaction-free-energy:exergonic}{exergónica} a \qty{298}{K}, aunque el término entrópico (\qty{+59.1}{kJ/mol})
juega en contra.
\end{solution}

\begin{solution}{exo:b2:reaction-free-energy:3}
$\Delta_r S^\circ = 69.95 - 130.68 - \tfrac12(205.15) = \qty{-163.3}{J/(K.mol)}$;
$\Delta_r G^\circ = -285.83 + 298.15 \times 0.1633 = \qty{-237.14}{kJ/mol}$.
Es la reacción de formación del agua líquida (un mol a partir de los elementos en
sus estados de referencia), así que su \omterm{def:b2:reaction-free-energy:reaction-gibbs}{energía de Gibbs estándar de reacción} es por
definición $\Delta_f G^\circ(\ce{H2O}, \text{l})$; las tablas dan el mismo
valor.
\end{solution}

\begin{solution}{exo:b2:reaction-free-energy:4}
Unos \qty{42}{J/(K.mol)} a \qty{298}{K} y \qty{87}{J/(K.mol)} a
\qty{1000}{K} (líquido). Saltos: \qty{10.6}{J/(K.mol)} a \qty{692.7}{K} y
\qty{97.7}{J/(K.mol)} a \qty{1180.2}{K}. $\Delta_{\text{fus}}H = 10.6 \times
692.7 = \qty{7.3}{kJ/mol}$ y $\Delta_{\text{vap}}H = 97.7 \times 1180.2 =
\qty{115}{kJ/mol}$.
\end{solution}

\begin{solution}{exo:b2:reaction-free-energy:5}
$\Delta_r H^\circ = \qty{44.00}{kJ/mol}$, $\Delta_r S^\circ = 188.84 - 69.95 =
\qty{118.9}{J/(K.mol)}$, $\Delta_r G^\circ = 44.00 - 298.15 \times 0.1189 =
\qty{+8.55}{kJ/mol}$. A \qty{298}{K} el cambio no es reversible (el agua líquida
y el vapor a \qty{1}{bar} no están en equilibrio), de modo que $\Delta S \ne
Q/T$: $\Delta_r H^\circ/T = \qty{147.6}{J/(K.mol)}$. El $\Delta_r G^\circ$
positivo significa que a \qty{298}{K} el agua no se evapora en
una atmósfera de vapor puro a \qty{1}{bar}: el vapor se condensa. Ambos
están en equilibrio donde $\Delta_r G^\circ = 0$: $T = 44\,004/118.9 =
\qty{370}{K}$, a \qty{3}{K} del punto de ebullición normal
(\qty{373}{K}): la \omterm{def:b2:reaction-free-energy:ellingham-approximation}{aproximación de Ellingham} es buena en \qty{75}{K}.
\end{solution}

\begin{solution}{exo:b2:reaction-free-energy:6}
$T_i = 57\,100/175.7 = \qty{325}{K}$. Por debajo, $\Delta_r G^\circ > 0$ y
se ve favorecido el \ce{N2O4} incoloro; enfriar el tubo desplaza el gas hacia
\ce{N2O4} y el color pardo del \ce{NO2} se desvanece (el vínculo cuantitativo con
la composición está en el \cref{ch:b2:equilibrium-shifts}).
\end{solution}

\begin{solution}{exo:b2:reaction-free-energy:7}
$\Delta_r S^\circ = -\dd\Delta_r G^\circ/\dd T = -(-4.0 + 12.0)/100 =
\qty{-0.080}{kJ/(K.mol)} = \qty{-80}{J/(K.mol)}$ and $\Delta_r H^\circ =
\Delta_r G^\circ + T\Delta_r S^\circ = -12.0 + 300 \times (-0.080) =
\qty{-36.0}{kJ/mol}$. Comprobación: $\Delta_r G^\circ/T$ pasa de $-0.0400$ a
$\qty{-0.0100}{kJ/(K.mol)}$, pendiente $3.00 \times 10^{-4}$ por kelvin, igual a
$-\Delta_r H^\circ/T^2 = 36.0/346^2 = 3.00 \times 10^{-4}$ a la temperatura
media geométrica $\sqrt{300 \times 400} = \qty{346}{K}$.
\end{solution}

\begin{solution}{exo:b2:reaction-free-energy:8}
$\Delta_r S^\circ = 186.25 - 5.74 - 2(130.68) = \qty{-80.8}{J/(K.mol)}$,
$\Delta_r G^\circ = -74.87 + 298.15 \times 0.0808 = \qty{-50.8}{kJ/mol}$,
el $\Delta_f G^\circ$ tabulado del metano. Al tener ambos términos el mismo
signo, el metano se vuelve inestable por encima de $T_i = 74\,870/80.8 \approx
\qty{926}{K}$ (\omterm{def:b2:reaction-free-energy:ellingham-approximation}{aproximación de Ellingham}): el metano caliente se craquea en carbono e
hidrógeno, una vía hacia ambos.
\end{solution}

\begin{solution}{exo:b2:reaction-free-energy:9}
Sustitución: \ce{CH3CHBrCH3 + OH- -> CH3CH(OH)CH3 + Br-}; eliminación:
\ce{CH3CHBrCH3 + OH- -> CH2=CHCH3 + H2O + Br-}. La diferencia de las dos
\omterm{def:b2:reaction-free-energy:reaction-gibbs}{energías de Gibbs estándar de reacción} es $\Delta(\Delta_r G^\circ) = 40 - T \times
0.090$, negativa por encima de $40\,000/90 \approx \qty{440}{K}$. La eliminación forma
una molécula más que la sustitución: su entropía es mayor, y el término
entrópico crece proporcionalmente a $T$.
\end{solution}

\begin{solution}{exo:b2:reaction-free-energy:10}
Cada una de las $N_A$ moléculas tiene 2 orientaciones: $W = 2^{N_A}$, luego $S =
k_B\ln 2^{N_A} = R\ln 2 = \qty{5.76}{J/(K.mol)}$. El cristal no está
perfectamente ordenado: a baja temperatura las moléculas quedan congeladas en sus
orientaciones aleatorias y no pueden alcanzar la disposición ordenada única que supone el tercer
principio.
\end{solution}

\begin{solution}{exo:b2:reaction-free-energy:11}
Se integra $\dd\Delta_r S^\circ/\dd T = c/T$: $\Delta_r S^\circ(T) = \Delta_r
S^\circ(T_0) + c\ln(T/T_0)$; Kirchhoff da $\Delta_r H^\circ(T) = \Delta_r
H^\circ(T_0) + c(T - T_0)$; se resta $T\Delta_r S^\circ(T)$. Para la caliza a
\qty{1100}{K}: Ellingham $178.32 - 1100 \times 0.16064 = \qty{1.62}{kJ/mol}$;
corrección $c[T - T_0 - T\ln(T/T_0)] = -0.00195 \times (801.85 - 1436.3) =
\qty{+1.24}{kJ/mol}$: $\Delta_r G^\circ = \qty{2.86}{kJ/mol}$. El error,
\qty{1.2}{kJ/mol}, es menos del uno por ciento de $\Delta_r H^\circ$, pero cerca de
la \omterm{def:b2:reaction-free-energy:inversion-temperature}{temperatura de inversión} decide el signo.
\end{solution}

\begin{solution}{exo:b2:reaction-free-energy:12}
(a) $\Delta_r H^\circ = -763.2 + 944.0 = \qty{+180.8}{kJ/mol}$, $\Delta_r
S^\circ = 353.2 + 205.15 - 50.62 - 2(223.08) = \qty{+61.6}{J/(K.mol)}$;
$\Delta_r G^\circ(1000) = 180.8 - 61.6 = \qty{+119.2}{kJ/mol}$: \omterm{def:b2:reaction-free-energy:exergonic}{endergónica},
no hay cloración. (b) Para \ce{C + O2 -> CO2}, $\Delta_r G^\circ(1000) =
-393.5 - 2.9 = \qty{-396.4}{kJ/mol}$. La suma, \ce{TiO2(s) + 2Cl2(g) + C(s)
-> TiCl4(g) + CO2(g)}, tiene $\Delta_r G^\circ = \qty{-277.2}{kJ/mol}$: el carbono
consume el dioxígeno formado por la primera reacción, y la reacción acoplada es
fuertemente \omterm{def:b2:reaction-free-energy:exergonic}{exergónica}.
\end{solution}

\begin{solution}{pb:b2:reaction-free-energy:1}
\textbf{1.} \ce{CaCO3(s) -> CaO(s) + CO2(g)}.
\textbf{2.} $\Delta_r H^\circ = -635.09 - 393.51 + 1206.92 =
\qty{+178.3}{kJ/mol}$: \omterm{def:b2:reaction-enthalpy:exothermic}{endotérmica}.
\textbf{3.} $\Delta_r S^\circ = 39.75 + 213.79 - 92.9 = \qty{+160.6}{J/(K.mol)}$,
positiva porque se forma un gas a partir de un sólido.
\textbf{4.} $\Delta_r G^\circ = 178.32 - 298.15 \times 0.16064 =
\qty{+130.4}{kJ/mol}$.
\textbf{5.} Sí: para la descomposición hace falta $\Delta_r G < 0$, y aquí es
fuertemente positiva; el dióxido de carbono del aire, muy por debajo de \qty{1}{bar},
rebaja $\Delta_r G$, pero mucho menos que \qty{130}{kJ/mol} (capítulo siguiente).
\textbf{6.} $\Delta_r C^\circ_p = 42.80 + 37.13 - 81.88 =
\qty{-1.95}{J/(K.mol)}$, minúscula: $\Delta_r H^\circ$ y $\Delta_r S^\circ$
apenas varían con $T$.
\textbf{7.} $\Delta_r G^\circ(T) = 178.32 - 0.16064\,T$ (\unit{kJ/mol}).
\textbf{8.} \qty{+17.7}{kJ/mol} a \qty{1000}{K}; \qty{-30.5}{kJ/mol} a
\qty{1300}{K}.
\textbf{9.} $T_i = 178\,320/160.64 = \qty{1110}{K}$.
\textbf{10.} $\Delta_r G^\circ/T = 178.32/T - 0.16064$, cuya derivada es
$-178.32/T^2 = -\Delta_r H^\circ/T^2$.
\textbf{11.} En $T_i$, $c[T_i - T_0 - T_i\ln(T_i/T_0)] = -0.00195 \times
(811.9 - 1459.4) = \qty{+1.26}{kJ/mol}$; $\Delta_r G^\circ$ se anula
$1260/160.6 \approx \qty{8}{K}$ más arriba, hacia \qty{1118}{K}: menos del
uno por ciento.
\textbf{12.} A \qty{1000}{K}, $\Delta_r G > 0$: la descomposición no puede
producirse (la cal absorbería dióxido de carbono). A \qty{1300}{K}, $\Delta_r G <
0$: la caliza se descompone.
\textbf{13.} $\Delta_r G$ se vuelve positiva: la cal capta el dióxido de carbono
y vuelve a ser carbonato.
\textbf{14.} $\delta S_{\text{c}} = -\Delta_r G\,\dd\xi/T = 30\,506/1300 =
\qty{23.5}{J/K}$ por mol.
\textbf{15.} Sí: el \cref{ch:b2:chemical-potential} mostrará que
$\Delta_r G = \Delta_r G^\circ + RT\ln(p_{\ce{CO2}}/p^\circ)$, que disminuye
cuando la presión del dióxido de carbono es inferior a \qty{1}{bar}, de modo que la
descomposición empieza por debajo de $T_i$.
\textbf{16.} $n = 10^6/56.08 = \qty{1.783e4}{mol}$.
\textbf{17.} $1.783 \times 10^4 \times 178.32 = \qty{3.18e6}{kJ}$, es decir,
\qty{3.18}{GJ} por tonelada.
\textbf{18.} $1.783 \times 10^4 \times 44.01 = \qty{785}{kg}$ de \ce{CO2}.
\textbf{19.} Calor que hay que aportar: $3.18 \times 10^6/0.50 = \qty{6.36e6}{kJ}$, luego
$6.36 \times 10^6/802.3 = \qty{7.93e3}{mol}$ de metano, \qty{127}{kg}.
\textbf{20.} $7.93 \times 10^3 \times 44.01 = \qty{349}{kg}$ de \ce{CO2}.
\textbf{21.} $785 + 349 = \qty{1134}{kg}$ por tonelada de cal, el 69\,\%
procedente de la caliza.
\textbf{22.} El dióxido de carbono de la caliza lo fija la estequiometría,
sea cual sea lo que calienta el horno; solo puede recortarse la parte del combustible.
\textbf{23.} La caliza se descompone bajo \qty{1}{bar} de dióxido de carbono por encima
de su \omterm{def:b2:reaction-free-energy:inversion-temperature}{temperatura de inversión}, $\boldsymbol{T_i \approx \qty{1.11e3}{K}}$
(unos \qty{840}{\celsius}).
\end{solution}
